\section*{Lampiran A: Source Code Firmware ESP32-S3 (C)}
\phantomsection
\addcontentsline{toc}{section}{Lampiran A: Source Code Firmware ESP32-S3}

Kode berikut merupakan firmware yang dijalankan pada ESP32-S3 di simulator
Wokwi. Program membaca suhu dari sensor DS18B20 dan mengirimkannya ke
ThingSpeak setiap 15 detik. Interval 15 detik dipilih sebagai batas aman
minimum \textit{update rate} ThingSpeak pada akun gratis.

\begin{lstlisting}[style=cstyle,
    caption={Firmware ESP32-S3 untuk Monitoring Suhu DS18B20 dan Pengiriman ke ThingSpeak},
    label={lst:esp32}]
#include <WiFi.h>
#include <ThingSpeak.h>
#include <OneWire.h>
#include <DallasTemperature.h>

// ==========================
// WiFi Wokwi
// ==========================
const char* ssid     = "Wokwi-GUEST";
const char* password = "";

// ==========================
// ThingSpeak
// ==========================
WiFiClient client;
unsigned long channelID     = 3413573;
const char*   writeAPIKey   = "ISI_WRITE_API_KEY_ANDA";

// ==========================
// DS18B20 (pin GPIO4)
// ==========================
#define ONE_WIRE_BUS 4
OneWire oneWire(ONE_WIRE_BUS);
DallasTemperature sensors(&oneWire);

// ==========================
// Interval kirim: 15 detik (batas aman ThingSpeak)
// ==========================
unsigned long lastSendTime  = 0;
const unsigned long sendInterval = 15000;

void connectWiFi() {
  Serial.print("Menghubungkan ke WiFi");
  WiFi.begin(ssid, password);
  while (WiFi.status() != WL_CONNECTED) {
    delay(500);
    Serial.print(".");
  }
  Serial.println();
  Serial.println("WiFi terhubung");
  Serial.print("IP Address: ");
  Serial.println(WiFi.localIP());
}

void checkWiFiConnection() {
  if (WiFi.status() != WL_CONNECTED) {
    Serial.println("WiFi terputus. Menghubungkan ulang...");
    WiFi.disconnect();
    WiFi.begin(ssid, password);
    while (WiFi.status() != WL_CONNECTED) {
      delay(500);
      Serial.print(".");
    }
    Serial.println();
    Serial.println("WiFi tersambung kembali");
  }
}

void setup() {
  Serial.begin(115200);
  sensors.begin();
  connectWiFi();
  ThingSpeak.begin(client);
  // Kirim langsung saat loop pertama
  lastSendTime = millis() - sendInterval;
  Serial.println("Sistem Monitoring Suhu DS18B20 Berbasis IoT Cerdas Dimulai");
  Serial.println("Interval pengiriman data ke ThingSpeak: 15 detik");
}

void loop() {
  checkWiFiConnection();

  sensors.requestTemperatures();
  float temperature = sensors.getTempCByIndex(0);

  if (temperature == DEVICE_DISCONNECTED_C) {
    Serial.println("Sensor DS18B20 tidak terbaca. Cek wiring.");
    delay(2000);
    return;
  }

  Serial.println("================================");
  Serial.print("Suhu aktual : ");
  Serial.print(temperature);
  Serial.println(" C");

  if (millis() - lastSendTime >= sendInterval) {
    // ESP32 hanya mengirim suhu aktual ke Field 1
    ThingSpeak.setField(1, temperature);
    int responseCode = ThingSpeak.writeFields(channelID, writeAPIKey);

    Serial.print("ThingSpeak response: ");
    Serial.println(responseCode);

    if (responseCode == 200) {
      Serial.println("Data suhu berhasil dikirim ke ThingSpeak Field 1");
    } else {
      Serial.println("Data suhu gagal dikirim ke ThingSpeak");
      Serial.println("Kemungkinan update terlalu cepat atau API Key salah.");
    }
    lastSendTime = millis();
  }

  delay(1000);
}
\end{lstlisting}
